\section{Practices for trademark-text research}\label{sec:practices}

Three practices follow for research that reads trademark text. First,
score distributions, not words. Word-level divergence is dominated by
description length --- on this corpus the correlation with log distinct
terms is $-0.64$ --- and firm-level performance correlations built
on it dissolve or reverse under distribution scoring on identical
samples. Any term-level novelty measure should be accompanied by the
theme-level rescoring that costs one model fit, and any firm-level
analysis should hold counsel of record, which the record observes on
every filing and which predicts reaching SEC reporting far more strongly
than any text measure. Second, treat the
registered description as a conformed document. When an application is
abandoned and the applicant tries again, the second attempt is rewritten
toward the class's typical language --- from both tails, and whoever
drafts it. Samples of registered marks therefore understate the
atypicality applicants first attempted, and estimates on registered text
are estimates on language that prosecution has already filtered. Studies using description
similarity or breadth inherit this compression. Third, the two components
of surprise are not interchangeable and should not be pooled. Lead is
robust to every partition, weighting, and window tried here;
atypicality's apparent effects can be artifacts of the partition (its
protective effect at the five-year proof belongs to theme importation,
visible only to global models) or of a control (its listing effect
depends on how description length is held). A study that reports
``novelty'' from trademark text without separating the temporal from the
cross-sectional component may attribute to one what belongs to the other.

For users of the classification itself, one caution: internet-bearing
filings in goods classes survive at rates closer to the technology
classes than to their own, so Nice classes used as industry controls do
not partition the internet economy. And for the maintenance data, one
asset: the five-year proof of continued use records, for every
registered mark, whether its owner was still willing to swear the goods
were being sold --- a product-level survival outcome with a date on it.
The indicator literature has largely left it unused, and it is
observable for any cohort more than seven years old at no cost beyond
parsing the prosecution history.
